% =============================================================================
% CARE Infrastructure Requirement Specification -- CARE Deployment
% -----------------------------------------------------------------------------
% Reusable, provider-agnostic template.
%
% HOW TO CUSTOMISE
%   Every deployment-specific value is a macro defined once in the
%   CONFIGURATION block below and referenced throughout the document.
%   1. DEFAULTS   -- every macro has a sensible default (\providecommand).
%                    Do NOT edit the defaults; the template stays reusable.
%   2. OVERRIDES  -- to change a value, copy its line into the OVERRIDES
%                    block (or into spec-overrides.tex) and change
%                    \providecommand to \newcommand. Overrides always win
%                    because they are defined first.
%   Values whose default renders as a red [placeholder] MUST be overridden
%   for a final document; everything else is optional.
%
% Compile with pdflatex (Overleaf default works fine).
% =============================================================================
\documentclass[11pt,a4paper]{article}

\usepackage[margin=2.5cm]{geometry}
\usepackage[T1]{fontenc}
\usepackage[utf8]{inputenc}
\usepackage{lmodern}
\usepackage{booktabs}
\usepackage{longtable}
\usepackage{array}
\usepackage{enumitem}
\usepackage{xcolor}
\usepackage{parskip}
\usepackage{lastpage}
\usepackage{fancyhdr}
\usepackage[hidelinks]{hyperref}

% Inline placeholder marker (red = value still needs to be provided)
\newcommand{\ph}[1]{\textcolor{red!70!black}{\texttt{[#1]}}}

% =============================================================================
% CONFIGURATION
% =============================================================================

% -----------------------------------------------------------------------------
% 1. OVERRIDES -- per-deployment values. Uncomment / add \newcommand lines
%    here, or keep them in a separate spec-overrides.tex next to this file.
% -----------------------------------------------------------------------------
\InputIfFileExists{spec-overrides.tex}{}{}

% \newcommand{\OrgName}{E-Gov}
% \newcommand{\ProjectName}{CARE HMIS}
% \newcommand{\DataCenterName}{Manipur State Data Center}
% \newcommand{\RegionLocation}{Manipur, India}
% \newcommand{\BaseDomain}{care.mn.gov.in}
% \newcommand{\DocDate}{20 August 2026}
% \newcommand{\DocOwner}{Jane Doe, jane@example.org}
% \newcommand{\KubeNodeCount}{10}

% -----------------------------------------------------------------------------
% 2. DEFAULTS -- one line per value. \providecommand is skipped when the
%    macro was already defined in the OVERRIDES block above.
% -----------------------------------------------------------------------------

% --- Document metadata -------------------------------------------------------
\providecommand{\OrgName}{\ph{ORGANISATION NAME}}
\providecommand{\ProjectName}{\ph{PROJECT NAME}}
\providecommand{\DataCenterName}{\ph{STATE DATA CENTER NAME}}
\providecommand{\DocVersion}{0.1 (Draft)}
\providecommand{\DocDate}{\today}
\providecommand{\DocOwner}{\ph{DOCUMENT OWNER / CONTACT}}
\providecommand{\DocAuthor}{\ph{Author}}

% --- Environment identity ----------------------------------------------------
\providecommand{\BaseDomain}{\ph{care.example.com}}   % wildcard DNS/TLS base domain
\providecommand{\EnvName}{\ph{staging / production}}
\providecommand{\RegionLocation}{\ph{region / location}}

% --- Kubernetes cluster ------------------------------------------------------
\providecommand{\KubeVersion}{1.35 or later}
\providecommand{\ControlPlaneSla}{99.99\%}
\providecommand{\KubeNodeCount}{7}    % integers only: totals are computed
\providecommand{\KubeNodeCpu}{8}      % vCPU per node (integer)
\providecommand{\KubeNodeRam}{16}     % GB per node (integer)
\providecommand{\KubeNodeDisk}{128 GB}

% --- Database cluster --------------------------------------------------------
\providecommand{\PgVersion}{17}
\providecommand{\DbNodeCount}{3}      % integers only: totals are computed
\providecommand{\DbNodeCpu}{8}        % vCPU per node (integer)
\providecommand{\DbNodeRam}{32}       % GB per node (integer)
\providecommand{\DbNodeDisk}{1 TB}
\providecommand{\ReadReplicaCount}{1}
\providecommand{\DbFailoverRto}{5 minutes}
\providecommand{\ReplicationLagThreshold}{30 seconds}
\providecommand{\DbBackupRetention}{30 days}
\providecommand{\DbBackupFrequency}{daily full + continuous WAL archiving}
\providecommand{\RestoreTestSchedule}{quarterly}
\providecommand{\DbMaintenanceWindow}{\ph{maintenance window}}

% --- Storage -----------------------------------------------------------------
\providecommand{\ObjectStorageCapacity}{1 TB}

% --- Security / compliance ---------------------------------------------------
\providecommand{\WafPolicy}{\ph{applicable policy}}
\providecommand{\RbacRoles}{cluster-admin, deployer, read-only auditor}
\providecommand{\AuditLogRetention}{365 days}
\providecommand{\ComplianceStandards}{CIS Kubernetes Benchmarks and applicable
  industry/state standards \ph{e.g.\ MeitY guidelines, ISO 27001}}
\providecommand{\FlowLogRetention}{90 days}

% --- Observability / DR ------------------------------------------------------
\providecommand{\AlertChannel}{\ph{email / chat / on-call tool}}
\providecommand{\ClusterRpo}{15 minutes}
\providecommand{\ClusterRto}{4 hours}
\providecommand{\DrDrillFrequency}{annually}

% --- CI/CD -------------------------------------------------------------------
\providecommand{\CicdTool}{\ph{GitHub Actions / GitLab CI / other}}
\providecommand{\ContainerRegistry}{\ph{to be assigned}}

% --- AI services -------------------------------------------------------------
\providecommand{\AiTextContext}{256k}
\providecommand{\AiTextOutput}{16k}
\providecommand{\AiMmluPro}{75\%}
\providecommand{\AiIfEval}{85\%}
\providecommand{\AiGpqa}{60\%}
\providecommand{\AiMmmu}{65\%}
\providecommand{\AiAudioMinutes}{30 minutes}
\providecommand{\AiWerEnglish}{15\%}
\providecommand{\AiWerHindi}{25\%}
\providecommand{\AiNoticeDays}{30 days}
\providecommand{\AiPromptTokens}{8{,}000}
\providecommand{\AiTtft}{2 seconds}
\providecommand{\AiTokensPerSecond}{30}
\providecommand{\AiRtf}{15}
\providecommand{\AiAvailability}{99.5\%}

% --- Service levels ----------------------------------------------------------
\providecommand{\PlatformSla}{99.9\% monthly}
\providecommand{\SupportHours}{\ph{24x7 / business hours}}
\providecommand{\IncidentResponse}{15 minutes}
\providecommand{\IncidentResolution}{4 hours}
\providecommand{\MaintenanceWindows}{\ph{schedule}}

% --- Environment values sheet ------------------------------------------------
\providecommand{\DicomEnabled}{\ph{yes / no}}
\providecommand{\IngressIp}{\ph{x.x.x.x}}
\providecommand{\EgressIp}{\ph{x.x.x.x}}
\providecommand{\SmtpHost}{\ph{smtp.example.org}}

% --- Derived values (computed; do not override) ------------------------------
\newcommand{\KubeTotalCpu}{\the\numexpr\KubeNodeCount*\KubeNodeCpu\relax}
\newcommand{\KubeTotalRam}{\the\numexpr\KubeNodeCount*\KubeNodeRam\relax~GB}
\newcommand{\DbTotalCpu}{\the\numexpr\DbNodeCount*\DbNodeCpu\relax}
\newcommand{\DbTotalRam}{\the\numexpr\DbNodeCount*\DbNodeRam\relax~GB}

% =============================================================================
% END CONFIGURATION -- no deployment-specific values below this line
% =============================================================================

% Requirement ID counter
\newcounter{reqctr}
\newcommand{\req}{\refstepcounter{reqctr}\textbf{REQ-\padreq}}
\makeatletter
\newcommand{\padreq}{\ifnum\value{reqctr}<10 0\fi\ifnum\value{reqctr}<100 0\fi\arabic{reqctr}}
\makeatother

% Header / footer
\pagestyle{fancy}
\fancyhf{}
\fancyhead[L]{\small \ProjectName}
\fancyhead[R]{\small CARE Infrastructure Requirement Specification}
\fancyfoot[L]{\small Version \DocVersion}
\fancyfoot[R]{\small Page \thepage\ of \pageref{LastPage}}
\renewcommand{\headrulewidth}{0.4pt}

\begin{document}

% -----------------------------------------------------------------------------
% Title page
% -----------------------------------------------------------------------------
\begin{titlepage}
  \centering
  \vspace*{3cm}
  {\LARGE\bfseries CARE Infrastructure Requirement Specification\par}
  \vspace{0.6cm}
  {\Large Hosting the CARE Software Platform\par}
  \vspace{0.3cm}
  {\large at \DataCenterName\par}
  \vspace{2cm}
  \begin{tabular}{ll}
    \toprule
    Organisation      & \OrgName \\
    Project           & \ProjectName \\
    Document Version  & \DocVersion \\
    Date              & \DocDate \\
    Document Owner    & \DocOwner \\
    \bottomrule
  \end{tabular}
  \vfill
  {\small This document is a reusable template. Values shown as
  \ph{placeholder} must be filled in for each deployment.\par}
\end{titlepage}

% -----------------------------------------------------------------------------
\section*{Revision History}
\begin{tabular}{@{}llll@{}}
  \toprule
  \textbf{Version} & \textbf{Date} & \textbf{Author} & \textbf{Description} \\
  \midrule
  \DocVersion & \DocDate & \DocAuthor & Initial draft \\
  \bottomrule
\end{tabular}

\tableofcontents
\newpage

% -----------------------------------------------------------------------------
\section{Purpose and Scope}

This document specifies the infrastructure required to host the
\textbf{CARE} software platform (an open-source healthcare information and
TeleICU platform) within \DataCenterName. It is written in
provider-agnostic terms: any infrastructure provider or on-premises platform that
satisfies the requirements herein (e.g.\ a CNCF-conformant Kubernetes
distribution, an S3-compatible object store, and a managed or self-hosted
PostgreSQL cluster) is acceptable.

The intended audience includes data center operations teams, infrastructure
vendors, system integrators, and the application maintenance team.

\subsection{Conventions}

Requirements are individually numbered (\textbf{REQ-nnn}) and use RFC~2119
keywords: \textbf{MUST}, \textbf{SHOULD}, \textbf{MAY}. Placeholder values are
shown as \ph{placeholder}.

% -----------------------------------------------------------------------------
\section{Compute -- High Availability Kubernetes Cluster}

\begin{itemize}[leftmargin=2.2cm]
  \item[\req] The platform MUST provide a dedicated CNCF-conformant
    Kubernetes cluster, version \KubeVersion, with a supported upgrade
    path. Only CARE-specific containers and supporting CARE platform services
    MUST run on this cluster.
  \item[\req] The control plane MUST be highly available: a multi-master
    configuration with at least \textbf{three} master/control-plane nodes (or
    an equivalent managed control plane with a published availability SLA of
    \ControlPlaneSla{} or better).
  \item[\req] Worker nodes MUST be distributed across multiple failure
    domains (availability zones, racks, or geo-locations) to avoid single
    points of failure.
  \item[\req] The cluster MUST support node auto-repair or documented
    procedures for replacing failed nodes without application downtime.
  \item[\req] A Cluster Autoscaler (or equivalent) MUST dynamically adjust
    the number of worker nodes based on resource utilisation, within the
    limits defined in Table~\ref{tab:sizing}.
\end{itemize}

\subsection{Required Capacity Sizing}

\begin{table}[h!]
  \centering
  \scriptsize
  \caption{Required infrastructure capacity.}
  \label{tab:sizing}
  \setlength{\tabcolsep}{4pt}
  \begin{tabular}{@{}lllrrrrr@{}}
    \toprule
    \textbf{Role} & \textbf{Type} & \textbf{Quantity} &
    \textbf{vCPU/VM} & \textbf{RAM/VM} & \textbf{Disk/VM} &
    \textbf{Total vCPU} & \textbf{Total RAM} \\
    \midrule
    Kubernetes workload & Virtual & \KubeNodeCount & \KubeNodeCpu & \KubeNodeRam~GB & \KubeNodeDisk & \KubeTotalCpu & \KubeTotalRam \\
    PostgreSQL cluster  & Virtual & \DbNodeCount   & \DbNodeCpu   & \DbNodeRam~GB   & \DbNodeDisk   & \DbTotalCpu   & \DbTotalRam   \\
    \bottomrule
  \end{tabular}
\end{table}

% -----------------------------------------------------------------------------
\section{Networking}

\subsection{Cluster Networking}
\begin{itemize}[leftmargin=2.2cm]
  \item[\req] The cluster MUST use a CNI (Container Network Interface) plugin
    that supports Kubernetes \texttt{NetworkPolicy} resources, enabling
    micro-segmentation between workloads.
  \item[\req] It MUST be possible to create Network Policy resources to
    control ingress and egress traffic between pods and namespaces.
  \item[\req] Dedicated, non-overlapping, adequately sized \texttt{/24} IP
    ranges MUST be allocated for nodes, pods, services, and databases.
  \item[\req] Outbound internet access from the cluster MUST be provided via
    NAT with a dedicated static egress IP address. The address MUST remain
    stable to support allow-listed API access to secure external endpoints.
\end{itemize}

\subsection{Ingress, DNS, and TLS}
\begin{itemize}[leftmargin=2.2cm]
  \item[\req] An Ingress controller or Kubernetes Gateway API implementation
    MUST be supported for L7 routing and traffic management, including SSL/TLS
    termination, header rewriting, and authentication hooks.
  \item[\req] Custom Resource Definitions (CRDs) MUST be permitted on the
    cluster; in particular, cert-manager CRDs (\texttt{ClusterIssuer},
    \texttt{Certificate}, etc.) for automated certificate issuance via
    Let's Encrypt / ACME or another certificate authority.
  \item[\req] Wildcard DNS for \texttt{*.\BaseDomain} MUST be provided
    so that CARE can assign subdomains for different platform purposes. A
    static public IP address MUST be reservable for the ingress endpoint.
  \item[\req] TLS certificates for \texttt{*.\BaseDomain} MUST be
    renewed automatically.
  \item[\req] A web application firewall (WAF) or edge security policy SHOULD
    be available, supporting OWASP core rule sets and geo-based access
    restrictions, per \WafPolicy.
\end{itemize}

% -----------------------------------------------------------------------------
\section{Object Storage (S3-Compatible)}

\begin{itemize}[leftmargin=2.2cm]
  \item[\req] An S3-compatible object storage service MUST be provided, with
    high availability and horizontal scalability.
  \item[\req] The object storage service MUST provide at least
    \ObjectStorageCapacity{} of initial capacity and MUST support expansion
    as required without service interruption.
  \item[\req] Data lifecycle policies MUST be configurable, covering object
    versioning, retention, and automatic deletion.
  \item[\req] Encryption MUST be enforced at rest (AES-256 or stronger,
    preferably with customer-managed keys) and in transit (TLS 1.2+).
  \item[\req] Fine-grained access control MUST be enforceable via bucket
    policies, IAM roles, and scoped access keys, so that only authorised
    users and applications can access stored data.
  \item[\req] HMAC-style access key/secret pairs (S3 API credentials) MUST be
    issuable per application, with the ability to rotate and revoke them.
  \item[\req] Presigned URL generation MUST be supported for time-limited
    client uploads and downloads.
\end{itemize}

% -----------------------------------------------------------------------------
\section{Persistent Storage (Kubernetes)}

\begin{itemize}[leftmargin=2.2cm]
  \item[\req] Dynamic volume provisioning MUST be supported via
    \texttt{StorageClass} resources backed by the data center's block or file
    storage, with the ability to define custom storage classes.
  \item[\req] SSD-class storage MUST be available for latency-sensitive
    workloads (databases, Redis persistence if enabled).
  \item[\req] Volume snapshots and data replication MUST be supported for
    critical application data, ensuring data integrity and availability.
  \item[\req] Volume expansion SHOULD be supported without workload
    recreation.
\end{itemize}

% -----------------------------------------------------------------------------
\section{Relational Database (PostgreSQL)}

\subsection{Cluster Topology and Availability}
\begin{itemize}[leftmargin=2.2cm]
  \item[\req] A managed, highly available PostgreSQL cluster (version
    \PgVersion{} or later) MUST be provided as the primary application
    datastore, with private-network-only connectivity and no public exposure.
    It MUST provide broad PostgreSQL extension support, including
    \texttt{pgvector}, supported Foreign Data Wrappers (FDWs), and other
    CARE-required extensions.
  \item[\req] The cluster MUST support one or more read replicas
    (\ReadReplicaCount{} initially) for scalability and fault tolerance.
  \item[\req] A separate PostgreSQL database or instance MUST be available
    for analytics workloads when required.
  \item[\req] Automated failover MUST be configured so that a standby is
    promoted with minimal downtime (target RTO \DbFailoverRto) upon primary
    failure.
  \item[\req] Replication lag MUST be monitored, with alerting when lag
    exceeds \ReplicationLagThreshold, enabling timely intervention.
\end{itemize}

\subsection{Data Protection and Security}
\begin{itemize}[leftmargin=2.2cm]
  \item[\req] Encryption MUST be enforced at rest (with managed key
    management / KMS integration) and in transit (TLS) for all database
    connections.
  \item[\req] Automated backups MUST be configured, combining full,
    incremental, and differential strategies, with point-in-time recovery
    (PITR). Retention: \DbBackupRetention; backup frequency:
    \DbBackupFrequency.
  \item[\req] Backup restoration MUST be tested on a defined schedule
    (\RestoreTestSchedule) with documented results.
  \item[\req] Database-level security policies MUST be enforced: role-based
    access control, audit logging of administrative actions, and least-
    privilege application accounts.
\end{itemize}

\subsection{Operations and Scaling}
\begin{itemize}[leftmargin=2.2cm]
  \item[\req] Routine maintenance jobs (index optimisation, vacuuming, data
    archiving) MUST be schedulable and monitored to sustain database
    performance.
  \item[\req] Data partitioning and sharding strategies SHOULD be supported
    where required to distribute load across nodes, with appropriate data
    modelling and management.
  \item[\req] Minor and major version upgrades MUST be plannable and
    executable with minimal downtime (target \DbMaintenanceWindow).
  \item[\req] Vertical and/or horizontal auto-scaling policies SHOULD be
    available to adjust database resources based on workload demand.
\end{itemize}

% -----------------------------------------------------------------------------
\section{Security}

\subsection{Access Control}
\begin{itemize}[leftmargin=2.2cm]
  \item[\req] Fine-grained Kubernetes RBAC MUST be enforced so that only
    authorised personnel can access and manage specific cluster resources.
    Roles: \RbacRoles.
  \item[\req] Pod-level security controls MUST be enforced via Pod Security
    Admission (or PodSecurityPolicies on older distributions), restricting
    privileged containers, host access, and capability escalation.
  \item[\req] Administrative access to nodes and databases MUST be restricted
    to a hardened bastion/jump host with SSH key-based authentication and
    session auditing.
\end{itemize}

\subsection{Secrets Management}
\begin{itemize}[leftmargin=2.2cm]
  \item[\req] A secrets management solution MUST be provided -- e.g.\
    HashiCorp Vault, a secret manager, or the Kubernetes Secrets Store
    CSI Driver -- for secure storage and distribution of sensitive data
    (database credentials, API keys, signing keys).
\end{itemize}

\subsection{Compliance and Auditing}
\begin{itemize}[leftmargin=2.2cm]
  \item[\req] Kubernetes audit logging MUST be enabled and retained for
    \AuditLogRetention.
  \item[\req] The deployment MUST maintain compliance with
    \ComplianceStandards.
  \item[\req] Network flow logs SHOULD be captured and retained for
    \FlowLogRetention{} for forensic analysis.
\end{itemize}

% -----------------------------------------------------------------------------
\section{Observability}

\begin{itemize}[leftmargin=2.2cm]
  \item[\req] Centralized logging MUST provide application error tracking,
    infrastructure and application log aggregation, searchable analysis, and
    appropriate retention controls. The specific technology stack will be
    selected during implementation.
  \item[\req] Metrics-based monitoring MUST track application and
    infrastructure performance (CPU, memory, disk, request latency, error
    rates, queue depth).
  \item[\req] Alerting MUST be configured for anomalies -- including failed
    deployments, pod crash loops, certificate expiry, database replication
    lag, storage capacity, and backup failures -- delivered to
    \AlertChannel.
  \item[\req] Dashboards SHOULD be provided for at-a-glance health of the
    cluster, databases, and application components.
\end{itemize}

% -----------------------------------------------------------------------------
\section{Backup and Disaster Recovery}

\begin{itemize}[leftmargin=2.2cm]
  \item[\req] A documented backup and disaster recovery strategy MUST exist,
    including off-site backups, data snapshots, and automated failover
    procedures.
  \item[\req] Full-cluster backup and restore procedures MUST be implemented,
    including etcd (or the managed control plane's state), to ensure
    integrity during failures.
  \item[\req] Recovery objectives MUST be agreed and tested:
    RPO \ClusterRpo, RTO \ClusterRto.
  \item[\req] DR drills MUST be conducted at least \DrDrillFrequency, with
    documented outcomes and remediation actions.
\end{itemize}

% -----------------------------------------------------------------------------
\section{Email (SMTP)}

\begin{itemize}[leftmargin=2.2cm]
  \item[\req] An SMTP server or relay service MUST be provided to handle
    email traffic for \BaseDomain{} (transactional email: password resets,
    notifications).
  \item[\req] SPF, DKIM, and DMARC records MUST be configured for the sending
    domain to ensure deliverability.
\end{itemize}

% -----------------------------------------------------------------------------
\section{CI/CD and Release Management}

\begin{itemize}[leftmargin=2.2cm]
  \item[\req] A container registry MUST be available, reachable from the
    cluster, for hosting application images.
  \item[\req] Continuous deployment MUST be supported so that successfully
    tested changes are promoted to production without manual intervention,
    using \CicdTool{} with short-lived federated credentials
    (OIDC / workload identity federation) rather than static keys.
  \item[\req] A centralized, access-controlled Git repository MUST serve as
    the source of truth for GitOps. It MUST include application deployment
    configuration, Kubernetes manifests, and infrastructure-as-code,
    including Terraform or equivalent definitions, with version history and
    change-approval controls.
\end{itemize}

% -----------------------------------------------------------------------------
\section{Documentation}

\begin{itemize}[leftmargin=2.2cm]
  \item[\req] Exhaustive documentation MUST be delivered, covering
    architecture, network topology, runbooks for routine operations, and
    escalation procedures.
  \item[\req] Documentation MUST be kept current with infrastructure changes
    as part of the change management process.
\end{itemize}

% -----------------------------------------------------------------------------
\section{Artificial Intelligence Services}

\begin{itemize}[leftmargin=2.2cm]
  \item[\req] An OpenAI-compatible HTTP API endpoint MUST be provided,
    implementing at minimum \texttt{/v1/chat/completions},
    \texttt{/v1/audio/transcriptions} and \texttt{/v1/embeddings}, with support
    for SSE streaming, tool calling, and JSON-Schema-constrained structured
    output. The capabilities below MAY be satisfied by separate models, but all
    MUST be reachable through this single endpoint.

  \item[\req] A text model MUST be available with a context window of at least
    \AiTextContext{} tokens and generated output of at least \AiTextOutput{}
    tokens, meeting or exceeding published scores of \AiMmluPro{} MMLU-Pro,
    \AiIfEval{} IFEval and \AiGpqa{} GPQA-Diamond.

  \item[\req] A vision-capable model MUST be available, accepting interleaved
    text and image input in a single request, meeting or exceeding a published
    score of \AiMmmu{} on MMMU (val). Document and form understanding
    MUST be supported at scanned-document resolution.

  \item[\req] A speech-transcription model MUST be available, accepting at
    least \AiAudioMinutes{} of audio per request, returning word-level
    timestamps, with automatic language identification and the ability to
    force a specified language. It MUST support code-mixed Hindi-English and
    Indian-English speech, and direct speech-to-English translation. Word
    Error Rate MUST NOT exceed \AiWerEnglish{} for English and \AiWerHindi{}
    for Hindi against a named public benchmark. For any other language
    offered, WER MUST be documented per language; multilingual averages are
    not acceptable in place of per-language figures.

  \item[\req] For each available model, the following MUST be documented and
    kept current: model family and version, parameter count, weight and
    KV-cache quantization format, serving runtime, and context window as
    served. Any change to these MUST be notified \AiNoticeDays{} in advance.
    Models MUST NOT be substituted, nor their precision, parameter count or
    context window reduced, without such notification.

  \item[\req] The service MUST sustain concurrent requests at a
    prompt length of \AiPromptTokens{} tokens, with p95 time-to-first-token
    not exceeding \AiTtft{} and sustained generation of at least
    \AiTokensPerSecond{} output tokens per second per stream. Speech
    transcription MUST achieve a real-time factor of at least \AiRtf,
    measured end-to-end at the API endpoint. Rate limits, where applied,
    MUST be documented numerically per model.

  \item[\req] Monthly availability MUST be at least \AiAvailability,
    measured at the API endpoint and independently verifiable. Rate limits,
    where applied, MUST be documented numerically per model.
\end{itemize}

% -----------------------------------------------------------------------------
\section{Service Levels and Support}

\begin{table}[h!]
  \centering
  \begin{tabular}{@{}ll@{}}
    \toprule
    \textbf{Item} & \textbf{Target} \\
    \midrule
    Platform availability            & \PlatformSla \\
    Support hours                    & \SupportHours \\
    Incident response (P1)           & \IncidentResponse \\
    Incident resolution (P1)         & \IncidentResolution \\
    Maintenance windows              & \MaintenanceWindows \\
    \bottomrule
  \end{tabular}
\end{table}

% -----------------------------------------------------------------------------
\appendix
\section{Glossary}

\begin{longtable}{@{}p{3.5cm}p{11cm}@{}}
  \toprule
  \textbf{Term} & \textbf{Definition} \\
  \midrule
  CNI    & Container Network Interface -- pluggable pod networking layer. \\
  CRD    & Custom Resource Definition -- extension mechanism for the Kubernetes API. \\
  HPA    & Horizontal Pod Autoscaler. \\
  KMS    & Key Management Service. \\
  PITR   & Point-In-Time Recovery. \\
  RBAC   & Role-Based Access Control. \\
  RPO/RTO & Recovery Point / Recovery Time Objective. \\
  WAF    & Web Application Firewall. \\
  WIF    & Workload Identity Federation -- keyless CI/CD authentication. \\
  \bottomrule
\end{longtable}

\section{Environment Values Sheet}

To be completed per environment (staging / production):

\begin{longtable}{@{}p{5.5cm}p{9cm}@{}}
  \toprule
  \textbf{Parameter} & \textbf{Value} \\
  \midrule
  Environment name              & \EnvName \\
  Region / location             & \RegionLocation \\
  Kubernetes version            & \KubeVersion \\
  Kubernetes workload nodes     & \KubeNodeCount{} virtual nodes; \KubeNodeCpu{} vCPU / \KubeNodeRam{} GB / \KubeNodeDisk{} each \\
  Pod CIDR                      & Dedicated non-overlapping \texttt{/24} \\
  Service CIDR                  & Dedicated non-overlapping \texttt{/24} \\
  Database version / size       & PostgreSQL \PgVersion{} or later; \DbNodeCount{} virtual nodes; \DbNodeCpu{} vCPU / \DbNodeRam{} GB / \DbNodeDisk{} each \\
  Read replica count            & \ReadReplicaCount \\
  Object storage capacity       & \ObjectStorageCapacity{} initially; scalable on demand \\
  Wildcard DNS domain           & \texttt{*.\BaseDomain} \\
  DICOM enabled                 & \DicomEnabled \\
  Ingress public IP             & \IngressIp \\
  Egress (NAT) IP               & \EgressIp \\
  SMTP relay host               & \SmtpHost \\
  Container registry            & \ContainerRegistry \\
  \bottomrule
\end{longtable}

\end{document}
